\reading{CR-4}{Capital Structure in Banks}{STRIKE FIRST}{10}

\begin{crkeybox}[Learning objectives for this reading]
\begin{enumerate}[label=\textbf{\alph*.},leftmargin=2em]
\item Evaluate a bank's economic capital relative to its level of credit risk.
\item Identify and describe important factors used to calculate economic capital for credit risk: probability of default, exposure, and loss rate.
\item Define and calculate expected loss (EL).
\item Define and calculate unexpected loss (UL).
\item Estimate the variance of default probability assuming a binomial distribution.
\item Calculate UL for a credit asset portfolio and the UL contribution of each asset under various scenarios of portfolio composition, asset characteristics and size.
\item Describe how economic capital is derived.
\item Explain how the credit loss distribution is modeled.
\item Describe challenges to quantifying credit risk.
\end{enumerate}
\end{crkeybox}

\begin{crtrapbox}[Notation warning for this reading]
This reading uses \term{EA} (exposure amount) and \term{LR} (loss rate) where
CR-3 used EAD and LGD. They are the same quantities. The source notes state
plainly that loss rate is \emph{also referred to as} loss given default, and that
$\mathrm{LR} = 1 - \mathrm{RR}$ where RR is the recovery rate. Both notations are
carried here as written rather than harmonised into one.
\end{crtrapbox}

% =====================================================================
\lo{CR-4 a}{Evaluate a bank's economic capital relative to its level of credit risk.}

\begin{crdefbox}[The three quantities, in order]
\term{Expected loss (EL)} is the best estimate of the devaluation of a risky
asset, representing the average potential loss the bank expects from its credit
portfolio.

\smallskip
\term{Unexpected loss (UL)} is what happens when events lead to much larger
losses than the expected ones. These can exceed the expected loss significantly.
To safeguard against unexpected losses and prevent impairment, banks need to hold
capital reserves.

\smallskip
\term{Economic capital} is the additional capital reserves that banks set aside
specifically to cover unexpected losses beyond their credit reserves. It acts as
a buffer to absorb large losses and maintain solvency.
\end{crdefbox}

% =====================================================================
\lo{CR-4 b}{Identify and describe important factors used to calculate economic capital for credit risk: probability of default, exposure, and loss rate.}

\begin{enumerate}
\item \term{Probability of default (PD).} The likelihood of a borrower defaulting
on their debt. A borrower may briefly default and then quickly correct the
situation by making a payment, or paying interest charges or penalties for missed
payments.
\item \term{Exposure amount (EA).} The dollar value at risk, stated as the
outstanding loan balance or a percentage of the credit line. In the source notes
this is described as the loan balance or the credit line maximum.
\item \term{Loss rate (LR)}, also referred to as \term{loss given default
(LGD)}. It represents the likely percentage loss in case of default. The severity
of default is crucial to creditors: a brief or low-loss default is less
concerning than a permanent or high-loss default. LR is related to the recovery
rate (RR), which reflects the potential for loss recovery.
\end{enumerate}

\begin{crfmlbox}[Loss rate and recovery rate]
\[ \mathrm{LR} \;=\; 1 - \mathrm{RR} \]
\vk{$\mathrm{RR}$ = recovery rate, the fraction of exposure recovered after
default. Both are percentages.}
\end{crfmlbox}

% =====================================================================
\lo{CR-4 c}{Define and calculate expected loss (EL).}

Expected loss is the anticipated decrease in the value of a risky asset held by a
bank: the average amount of money a bank anticipates losing on a risky asset,
like a loan, that it holds on its balance sheet. EL represents the average
behaviour of the risky asset, but it does not capture the variation, or
unexpected loss.

\begin{crfmlbox}[Expected loss]
\[ \mathrm{EL} \;=\; \mathrm{EA} \times \mathrm{PD} \times \mathrm{LR} \]
\vk{$\mathrm{EA}$ = exposure amount (currency); $\mathrm{PD}$ = probability of
default (percentage); $\mathrm{LR}$ = loss rate (percentage). The source notes
also write this as $\mathrm{EL} = \mathrm{EAD} \times \mathrm{PD} \times
\mathrm{LGD}$, which is the same equation.}
\end{crfmlbox}

Unexpected loss arises from unanticipated default or credit migration. The
difference between the actual loss and the expected loss is called the
unexpected loss (UL). The sum $\mathrm{EL} + \mathrm{UL}$ is the total potential
loss, which is what matters for risk management.

\begin{itemize}
\item \term{Default probability} can be estimated using historical data or models
like the Merton model, which views the firm as holding a call option with a
strike price equal to its outstanding debt value.
\item \term{Credit migration} denotes the possible deterioration in
creditworthiness of the borrower. While a shift in migration may not result in
immediate default, the probability of such an event increases. It is also
possible for the reverse to occur, that is, the credit quality of the obligor
improves over time.
\end{itemize}

% =====================================================================
\lo{CR-4 d}{Define and calculate unexpected loss (UL).}

The observation that the unexpected loss represents the variability of potential
losses can be modelled using the typical definition of standard deviation:

\[ \mathrm{UL}_H \;\equiv\; \sqrt{\operatorname{var}(V_H)} \]

In the following equation the subscript $H$ is dropped, but note that the focus
throughout is on the horizon date $H$. After some algebra we derive the following
expression.

\begin{crfmlbox}[Unexpected loss for a single asset]
\[ \mathrm{UL} \;=\; \mathrm{EA} \times \sqrt{\mathrm{PD}\times\sigma^{2}_{\mathrm{LR}}
   \;+\; \mathrm{LR}^{2}\times\sigma^{2}_{\mathrm{PD}}} \]
\vk{$\sigma_{\mathrm{LR}}$ = standard deviation of the loss rate;
$\sigma_{\mathrm{PD}}$ = standard deviation of the probability of default. Note
the asymmetry inside the root: PD enters the first term \emph{unsquared} while LR
enters the second term \emph{squared}.}
\end{crfmlbox}

Lower default uncertainty and lower recovery uncertainty give a lower UL.

% =====================================================================
\lo{CR-4 e}{Estimate the variance of default probability assuming a binomial distribution.}

Default is a binary event: the borrower either defaults or does not. Treating it
as a Bernoulli trial gives the variance of the default probability directly.

\begin{crfmlbox}[Variance of the default probability]
\[ \sigma^{2}_{\mathrm{PD}} \;=\; \mathrm{PD} \times (1-\mathrm{PD}) \]
\vk{This is the binomial (Bernoulli) variance. It is the quantity that feeds the
second term under the root in the UL formula above.}
\end{crfmlbox}

\begin{crexbox}[XYZ bank: expected and unexpected loss]
Suppose XYZ bank has booked a loan with the following characteristics: total
commitment of \$2,000,000 of which \$1,800,000 is currently outstanding. The bank
has assessed an internal credit rating equivalent to a 1\% default probability
over the next year. The bank has additionally estimated a 40\% loss rate if the
borrower defaults. The standard deviation of PD and LR is 10\% and 30\%
respectively.

Calculate the expected and unexpected loss for XYZ bank.
\tcblower
\textbf{Expected loss.}
\[ \mathrm{EL} = \$1{,}800{,}000 \times 0.01 \times 0.40 = \$7{,}200 \]

\textbf{Unexpected loss.}
\[ \mathrm{UL} = \mathrm{EA}\times\sqrt{\mathrm{PD}\times\sigma^{2}_{\mathrm{LR}}
   + \mathrm{LR}^{2}\times\sigma^{2}_{\mathrm{PD}}} \]
\[ \mathrm{UL} = \$1{,}800{,}000 \times \sqrt{0.01\times 0.3^{2} + 0.4^{2}\times 0.1^{2}}
   \;=\; \$90{,}000 \]

Step by step: $0.01\times0.09 = 0.0009$; $0.16\times0.01 = 0.0016$; the sum is
$0.0025$; its square root is $0.05$; and $\$1{,}800{,}000 \times 0.05 = \$90{,}000$.
\end{crexbox}

\begin{crtrapbox}[Which exposure figure is used]
The exposure amount in both calculations is the \$1,800,000 \emph{outstanding},
not the \$2,000,000 total commitment. The commitment figure plays no part in the
arithmetic.
\end{crtrapbox}

% =====================================================================
\lo{CR-4 f}{Calculate UL for a credit asset portfolio and the UL contribution of each asset under various scenarios of portfolio composition, asset characteristics and size.}

Expected loss on the portfolio, $\mathrm{EL}_P$, is simply the sum of the
expected losses of each asset. Portfolio unexpected loss, $\mathrm{UL}_P$, is
more complicated because of the cross-terms in the variance formula for an
$N$-asset portfolio.

\begin{crfmlbox}[Portfolio expected and unexpected loss]
\[ \mathrm{EL}_P \;=\; \sum_i \mathrm{EL}_i \;=\; \sum_i \bigl(\mathrm{EA}_i \times
   \mathrm{LR}_i \times \mathrm{PD}_i\bigr) \]
\[ \mathrm{UL}_P \;=\; \sqrt{\;\sum_i\sum_j \rho_{ij}\,\mathrm{UL}_i\,\mathrm{UL}_j\;} \]
\vk{$\rho_{ij}$ = default correlation between assets $i$ and $j$. In the special
case where every $\rho_{ij}=1$ for $i \neq j$, $\mathrm{UL}_P$ collapses to the
simple sum of the individual unexpected losses.}
\end{crfmlbox}

\subsection*{Risk contribution (unexpected loss contribution)}

The \term{risk contribution (RC)}, also known as the \term{unexpected loss
contribution (ULC)}, measures how much of the portfolio's UL is attributable to a
single asset.

\gapbadge

{\footnotesize\color{FRMGrey}The source notes state the definition of RC and then
leave the general formula blank, giving only the two-asset case and the large-$n$
shortcut. The general expression below is written from the GARP curriculum.}

\begin{crgapbox}[What the objective asks for]
The objective asks for the UL contribution of each asset \emph{under various
scenarios of portfolio composition, asset characteristics and size}. The two-asset
case and the large-portfolio shortcut are both special cases of one general
formula, and the derivation chain between them is what makes the shortcut usable.
\end{crgapbox}

\begin{crfmlbox}[General risk contribution]
\[ \mathrm{ULC}_i \;=\; \mathrm{ULMC}_i \times \mathrm{UL}_i
   \;=\; \frac{\displaystyle\sum_{j} \mathrm{UL}_j\,\rho_{ij}}{\mathrm{UL}_P}\times \mathrm{UL}_i \]
\vk{$\mathrm{ULMC}_i$ = the marginal contribution of loan $i$ to portfolio
unexpected loss, obtained as the first partial derivative of $\mathrm{UL}_P$ with
respect to $\mathrm{UL}_i$. The formula has the important property that the sum of
the ULCs of all loans equals the portfolio-level UL.}
\end{crfmlbox}

For a \term{two-asset portfolio}, the risk contributions of each asset are:

\begin{crfmlbox}[Two-asset risk contributions]
\[ \mathrm{RC}_1 = \mathrm{UL}_1 \times
   \frac{\mathrm{UL}_1 + (\rho_{1,2}\times\mathrm{UL}_2)}{\mathrm{UL}_P}
   \qquad
   \mathrm{RC}_2 = \mathrm{UL}_2 \times
   \frac{\mathrm{UL}_2 + (\rho_{1,2}\times\mathrm{UL}_1)}{\mathrm{UL}_P} \]
\[ \mathrm{RC}_1 + \mathrm{RC}_2 = \mathrm{UL}_P \]
\end{crfmlbox}

For a \term{large portfolio of similar credits} with a constant pairwise
correlation $\rho$, the same expression reduces. Writing $\mathrm{UL}_P =
\mathrm{UL}_i\sqrt{n + \rho(n^{2}-n)}$ for $n$ identical loans and dividing by
$n$ gives
\[ \mathrm{ULC}_i \;=\; \frac{\mathrm{UL}_P}{n}
   \;=\; \mathrm{UL}_i\sqrt{\tfrac{1}{n} + \rho\bigl(1-\tfrac{1}{n}\bigr)}\,,
   \qquad\text{which for large } n \text{ reduces to} \]

\begin{crfmlbox}[Large-portfolio shortcut]
\[ \mathrm{ULC}_i \;=\; \mathrm{UL}_i \times \sqrt{\rho} \]
\vk{$\rho$ = the constant pairwise default correlation. This holds only when the
loans have approximately the same characteristics and size and $n$ is large.}
\end{crfmlbox}

\begin{crexbox}[Bigger Bank: portfolio EL, UL and risk contributions]
Bigger Bank has two assets outstanding. Assuming a correlation of 0.3 between the
assets, compute $\mathrm{EL}_P$ and $\mathrm{UL}_P$ as well as the risk
contribution of each asset.

\smallskip
\centering\footnotesize
\begin{tabular}{l r r}
\toprule
 & \thead{Asset A} & \thead{Asset B} \\
\midrule
EA & \$8,250,000 & \$1,800,000 \\
PD & 0.50\% & 1.00\% \\
LR & 50.00\% & 40.00\% \\
$\sigma_{\mathrm{PD}}$ & 7.05\% & 10.00\% \\
$\sigma_{\mathrm{LR}}$ & 25.00\% & 30.00\% \\
\bottomrule
\end{tabular}
\tcblower
\textbf{Step 1. Compute EL for both assets.}
\begin{align*}
\mathrm{EL}_A &= \mathrm{EA}\times\mathrm{PD}\times\mathrm{LR}
 = \$8{,}250{,}000 \times 0.005 \times 0.50 = \$20{,}625 \\
\mathrm{EL}_B &= \$1{,}800{,}000 \times 0.01 \times 0.40 = \$7{,}200
\end{align*}

\textbf{Step 2. Compute UL for both assets.}
\begin{align*}
\mathrm{UL}_A &= \$8{,}250{,}000 \times \sqrt{0.005\times0.25^{2} + 0.5^{2}\times0.0705^{2}}
 = \$325{,}333 \\
\mathrm{UL}_B &= \$1{,}800{,}000 \times \sqrt{0.01\times0.3^{2} + 0.4^{2}\times0.1^{2}}
 = \$90{,}000
\end{align*}

\textbf{Step 3. Compute $\mathrm{EL}_P$.}
\[ \mathrm{EL}_P = \$20{,}625 + \$7{,}200 = \$27{,}825 \]

\textbf{Step 4. Compute $\mathrm{UL}_P$.}
\[ \mathrm{UL}_P = \sqrt{(325{,}333)^{2} + (90{,}000)^{2} + 2(0.3)(325{,}333)(90{,}000)}
 = \$362{,}642 \]

\textbf{Step 5. Compute RC for both assets.}
\begin{align*}
\mathrm{RC}_A &= (325{,}333)\,\frac{325{,}333 + 0.3\times90{,}000}{362{,}642} = \$316{,}084 \\
\mathrm{RC}_B &= (90{,}000)\,\frac{90{,}000 + 0.3\times325{,}333}{362{,}642} = \$46{,}558 \\
\mathrm{RC}_A + \mathrm{RC}_B &= \$362{,}642 = \mathrm{UL}_P \quad\checkmark
\end{align*}
\end{crexbox}

\begin{crexbox}[The same portfolio at a lower correlation]
Using the information provided in the previous example and assuming the
correlation between the assets has decreased to 0.1, compute $\mathrm{EL}_P$ and
$\mathrm{UL}_P$.
\tcblower
\[ \mathrm{EL}_P = \$27{,}825 \]
\[ \mathrm{UL}_P = \sqrt{(325{,}333)^{2} + (90{,}000)^{2} + 2(0.1)(325{,}333)(90{,}000)}
 = \$346{,}118 \]
Because correlation does not impact each asset individually, the expected loss on
the portfolio remains the same. However, the unexpected loss (variation) has
decreased to \$346,118.
\end{crexbox}

\begin{center}
\begin{tikzpicture}
\begin{axis}[
  width=12.6cm, height=6.2cm,
  xlabel={Pairwise default correlation $\rho$},
  ylabel={Portfolio loss (\$000)},
  xlabel style={font=\sffamily\scriptsize}, ylabel style={font=\sffamily\scriptsize},
  tick label style={font=\sffamily\scriptsize},
  xmin=0, xmax=1, ymin=0, ymax=470,
  ytick={0,100,200,300,400},
  grid=major, grid style={FRMRule!50, line width=0.3pt},
  legend style={font=\scriptsize\sffamily, draw=FRMRule, fill=white,
    at={(0.5,-0.30)}, anchor=north, legend columns=-1, column sep=10pt}]
\addplot[domain=0:1, samples=120, smooth, draw=FRMNavy, line width=1.1pt]
  {sqrt(325.333^2 + 90^2 + 2*x*325.333*90)};
\addlegendentry{$\mathrm{UL}_P$}
\addplot[domain=0:1, samples=2, draw=FRMGreen, line width=1.1pt, sharp plot]
  {27.825};
\addlegendentry{$\mathrm{EL}_P$}
\addplot[only marks, mark=*, mark size=1.7pt, draw=FRMRed, fill=FRMRed]
  coordinates {(0.1,346.118) (0.3,362.642) (1,415.333)};
\addlegendentry{worked values}
\node[font=\sffamily\scriptsize, fill=white, inner sep=1.5pt, text=FRMRed]
  at (axis cs:0.30,300) {346.1 at $\rho=0.1$};
\draw[-{Stealth[length=4pt]}, FRMRed] (axis cs:0.20,312) -- (axis cs:0.11,338);
\node[font=\sffamily\scriptsize, fill=white, inner sep=1.5pt, text=FRMRed]
  at (axis cs:0.63,318) {415.3 at $\rho=1$};
\draw[-{Stealth[length=4pt]}, FRMRed] (axis cs:0.80,325) -- (axis cs:0.97,405);
\end{axis}
\end{tikzpicture}
\end{center}
\figcap{Both curves are computed from the Bigger Bank figures using the
portfolio formulas above. Correlation moves $\mathrm{UL}_P$ and leaves
$\mathrm{EL}_P$ completely flat, which is the point the worked pair of examples
is making. Note also how shallow the curve is: taking correlation all the way
from 0 to 1 moves $\mathrm{UL}_P$ only from 337.6 to 415.3, because the square
root compresses it.}

\subsection*{Diversifiable and undiversifiable risk}

When an asset is held in isolation it carries both diversifiable
(firm-specific) and undiversifiable (market) risk. Diversifiable risk can be
reduced to nearly zero in a large portfolio as asset-specific risks offset each
other. However, undiversifiable risk remains as the residual credit risk that
cannot be eliminated through diversification.

\subsection*{The effect of correlation}

The correlation between bank assets is crucial in assessing potential portfolio
losses. Higher correlation indicates concentration risk, where default in one
asset can affect others.

Estimating default correlations between unrelated obligors is challenging, and
the sheer number of calculations needed is unmanageable.

\begin{crkeybox}[The one-line version]
High correlation $=$ higher concentration risk.
\end{crkeybox}

% =====================================================================
\lo{CR-4 g}{Describe how economic capital is derived.}

The amount of economic capital needed depends on the distance between the
unexpected outcome and the expected outcome at a certain confidence level (often
99.97\%). This difference is estimated based on the shape of the loss
distribution, including the expected loss (EL) and the unexpected loss (UL).

\term{Capital multiplier (CM).} To find the economic capital for the portfolio,
the difference between the unexpected outcome and the confidence level is
represented as a multiple of the portfolio's unexpected loss.

\begin{crfmlbox}[Economic capital]
\[ \text{Economic Capital}_P \;=\; \mathrm{UL}_P \times \mathrm{CM}
   \qquad\qquad
   \text{Economic Capital}_i \;=\; \mathrm{ULC}_i \times \mathrm{CM} \]
\vk{$\mathrm{CM}$ = capital multiplier, the number of portfolio ULs between the
expected outcome and the chosen confidence level. Because the sum of the
$\mathrm{ULC}_i$ equals $\mathrm{UL}_P$, economic capital at the single
transaction level is directly proportional to that transaction's contribution to
overall portfolio credit risk.}
\end{crfmlbox}

\begin{center}
\begin{tikzpicture}
\begin{axis}[
  width=13.0cm, height=5.8cm,
  axis lines=left, xlabel={Loss over one year}, ylabel={Probability},
  xlabel style={font=\sffamily\scriptsize}, ylabel style={font=\sffamily\scriptsize},
  xtick=\empty, ytick=\empty,
  xmin=0, xmax=8, ymin=0, ymax=0.30, clip=true]
\addplot[domain=0:8, samples=220, smooth, draw=FRMNavy, line width=1pt] {0.55*x^1.4*exp(-1.15*x)};
\draw[FRMGreen, line width=0.9pt] (axis cs:2.09,0) -- (axis cs:2.09,0.20);
\node[above, font=\sffamily\scriptsize\bfseries, text=FRMGreen, fill=white,
  inner sep=1.5pt] at (axis cs:2.09,0.20) {$\mathrm{EL}_P$};
\draw[FRMRed, line width=0.9pt] (axis cs:5.6,0) -- (axis cs:5.6,0.115);
\node[above, font=\sffamily\scriptsize\bfseries, text=FRMRed, fill=white,
  inner sep=1.5pt] at (axis cs:5.6,0.115) {confidence level (e.g.\ 99.97\%)};
\draw[{Stealth[length=5pt]}-{Stealth[length=5pt]}, FRMGold, line width=0.9pt]
  (axis cs:2.09,0.055) -- (axis cs:5.6,0.055);
\node[font=\sffamily\scriptsize\bfseries, text=FRMGold, fill=white, inner sep=2pt]
  at (axis cs:3.85,0.055) {economic capital $=\mathrm{CM}\times\mathrm{UL}_P$};
\end{axis}
\end{tikzpicture}
\end{center}
\figcap{Illustrative shape only. Economic capital is the \emph{distance} between
the expected outcome and the chosen confidence level, expressed as a multiple of
$\mathrm{UL}_P$. Raising the confidence level pushes the right-hand line out and
raises the capital multiplier, which raises economic capital; it does not change
$\mathrm{EL}_P$.}

\begin{crexbox}[Two banks, same everything except correlation]
A risk analyst at a credit ratings agency is evaluating the economic capital for
credit risk of two competing regional banks, Bank ABC and Bank XYZ. The two banks
have the same credit asset exposure, duration of credit exposure, credit rating,
and expected loss. Assuming the average pairwise default correlation between
credit assets of Bank ABC is lower than that of Bank XYZ, and the two banks assess
their risk appetite at the same predetermined confidence level, which of the
following statements would be correct?

\smallskip
A. If the confidence level for both banks is increased, the level of economic
capital needed for Bank ABC and for Bank XYZ will both increase.

B. If the confidence level for both banks is decreased, the level of economic
capital needed will increase for Bank ABC but will decrease for Bank XYZ.

C. If the confidence level for both banks is increased, the level of economic
capital needed will decrease for Bank ABC but will increase for Bank XYZ.

D. If the confidence level for both banks is kept unchanged, the level of
economic capital needed for Bank ABC and for Bank XYZ will be equal.
\tcblower
\textbf{A is correct.} If the confidence level is raised, economic capital needed
would be higher since a higher capital multiplier is applied to the UL.

\textbf{B is incorrect.} A lower confidence level implies a lower credit rating
target (and a lower capital multiplier), and hence a lower capital level is
needed to cover credit assets, for \emph{both} banks.
\end{crexbox}

\begin{crexbox}[UL contribution in a large homogeneous portfolio]
A credit risk manager at a bank is estimating the unexpected loss (UL) of the
bank's portfolio of loans and the UL contributions of the individual loans to the
overall portfolio UL. The portfolio consists of a large number of loans and the
manager assumes that the loans have approximately the same characteristics and
size, with a constant pairwise default correlation of 0.32. Assuming the UL of
each loan is USD 10,500, what is the approximate UL contribution of each loan to
the portfolio UL?

\smallskip
A. USD 0 \qquad B. USD 1,075 \qquad C. USD 3,360 \qquad D. USD 5,940
\tcblower
\textbf{D is correct.} For a large number of loans, the unexpected loss
contribution of each loan to the portfolio UL is derived as:
\[ \mathrm{ULC}_i = \mathrm{UL}_i \times \sqrt{\rho}
 = 10{,}500 \times \sqrt{0.32} = \mathrm{USD}\ 5{,}939.70 \]
where $\mathrm{UL}_i = \mathrm{USD}\ 10{,}500$ and $\rho = 0.32$.
\end{crexbox}

% =====================================================================
\lo{CR-4 h}{Explain how the credit loss distribution is modeled.}

When estimating economic capital for credit risk, we focus on the tail of the
loss distribution because credit risks are not normally distributed and tend to
be highly skewed.

\begin{itemize}
\item To model credit events, a \term{beta distribution} is commonly used as it
allows losses to be defined between 0\% and 100\%. This distribution is flexible
and can be symmetric or skewed based on its shape parameters, which determine
symmetry and tail behaviour.
\item Modelling the tail of the credit loss distribution is more challenging and
often involves combining the beta distribution with a \term{Monte Carlo
simulation} for accuracy.
\end{itemize}

% =====================================================================
\lo{CR-4 i}{Describe challenges to quantifying credit risk.}

The bottom-up credit risk measurement framework has several limitations.

\begin{enumerate}
\item \term{Illiquidity.} It assumes credits are illiquid assets, so credit
losses are measured without considering correlations among risk factors as in
liquid markets. Illiquidity of credit assets poses challenges in accurately
quantifying credit risk.
\item \term{Short horizon.} Credit risk models often use only a one-year
estimation horizon, making it challenging to incorporate long-term changes in
borrower credit quality rather than capturing longer-term credit quality changes.
\item \term{Siloed risk types.} Other risks like operational and market risk are
separated and managed differently within the bank, not fully integrated with
credit risk assessment, leading to fragmented risk management.
\end{enumerate}
